% !TeX root = ../../tesis.tex
%%%%%%%%%%%%%%%%%%%%%%%%%%%%%%%%%%%%%%%%%%%%%%%%%%%%%%%%%%%%%%%%%%%
% SECCIÓN 5.3.2 — NIVEL DE COBERTURA ESPACIAL (C)
%%%%%%%%%%%%%%%%%%%%%%%%%%%%%%%%%%%%%%%%%%%%%%%%%%%%%%%%%%%%%%%%%%%
\subsection{Nivel de Cobertura Espacial (\texorpdfstring{$C$}{C})}
\label{subsec:cobertura-espacial}

El \textbf{Nivel de Cobertura Espacial} ($C$) evalúa la equidad territorial
de la red al cuantificar qué proporción del área de cada demarcación
político-administrativa se encuentra dentro de un radio de caminata aceptable
respecto a alguna estación. Se define como:

\begin{equation}
    C = \frac{A_{\text{cubierta}}}{A_{\text{total}}} \times 100
    \label{eq:cobertura-res}
\end{equation}

donde $A_{\text{cubierta}}$ es el área resultante de la unión de los
isócronos peatonales (\termino{buffers} de 800~m) de todas las estaciones del
sistema sobre la demarcación, y $A_{\text{total}}$ es la superficie geodésica
total de la demarcación analizada. El análisis abarca las 16~alcaldías de la
\cdmx{} y los 125 municipios de la \zmvm{} incluidos en el dataset
del \textbf{Apímetro} \citep{cabrera2025apimetro}, para un total de
141 demarcaciones.

La Figura~\ref{fig:grafica-cobertura-sistema} sintetiza el nivel de cobertura
agregado por sistema, evidenciando la marcada diferencia entre la red capilar
de superficie y los sistemas estructurados de mayor escala.

\begin{figure}[H]
    \centering
    \includegraphics[width=0.92\textwidth]{%
        Figures/Mapas/Tableau_Resources/RedActual/VFTModel/UniGrafica/%
        Grafica_VFTModel_2026_Cobertura_sistema}
    \caption[Cobertura territorial agregada por sistema de transporte]{
    Cobertura territorial agregada por sistema de transporte en la \zmvm{},
    calculada a partir de isócronos peatonales de 800~m sobre las
    141 demarcaciones analizadas. Corrida: 2026-04-18.}
    \label{fig:grafica-cobertura-sistema}
    \fuentefigura{Elaboración propia con \textbf{Tableau};
    datos: \texttt{VFTModel} \citep{cabrera2026VFTModel}.}
\end{figure}

% ---------------------------------------------------------------
\subsubsection{Sistemas de Infraestructura Estructurada}
\label{subsubsec:cobertura-estructurados}
% ---------------------------------------------------------------

Los sistemas con derecho de vía exclusivo o confinado —\textsc{Metro},
Metrobús, Tren Ligero, Cablebús, Tren Suburbano, Mexicable y Mexibús— presentan
coberturas territoriales altas en las alcaldías centrales, pero decrecen
abruptamente hacia la periferia. Su distribución refleja la lógica de inversión
en infraestructura pesada, concentrada históricamente en las zonas de mayor
densidad y actividad económica de la cuenca.

\begin{table}[H]
    \centering
    \caption[Cobertura territorial — sistemas estructurados, top 3 por alcaldía]{
    Tres alcaldías de la \cdmx{} con mayor cobertura territorial por sistema de
    infraestructura estructurada. Isócrono: 800~m. Corrida: 2026-04-18.}
    \label{tab:cobertura-estructurados}
    \fontsize{8.5pt}{10.5pt}\selectfont
    \begin{tabularx}{\textwidth}{|l|X|r|r|r|}
        \hline
        \textbf{Sistema} & \textbf{Alcaldía} &
        \textbf{Área total (km²)} & \textbf{Área cub. (km²)} &
        \textbf{$C$ (\%)} \\
        \hline
        \multirow{3}{*}{\textsc{Metro}}
            & Cuauhtémoc          & 32.50 & 28.35 & 87.23 \\
            & Benito Juárez       & 26.68 & 21.34 & 79.99 \\
            & Venustiano Carranza & 33.84 & 24.39 & 72.77 \\
        \hline
        \multirow{3}{*}{Metrobús}
            & Cuauhtémoc    & 32.50 & 26.74 & 82.27 \\
            & Benito Juárez & 26.68 & 16.21 & 60.74 \\
            & Iztacalco     & 23.08 & 13.95 & 60.44 \\
        \hline
        \multirow{3}{*}{Tren Ligero}
            & Coyoacán   &  53.88 &  9.31 & 17.28 \\
            & Xochimilco & 114.03 &  6.82 &  5.98 \\
            & Tlalpan    & 314.25 &  4.07 &  1.30 \\
        \hline
        \multirow{3}{*}{Cablebús}
            & Gustavo A. Madero &  87.84 & 11.31 & 12.88 \\
            & Miguel Hidalgo    &  46.39 &  5.50 & 11.85 \\
            & Iztapalapa        & 113.07 & 12.74 & 11.27 \\
        \hline
        \multirow{3}{*}{Tren Suburbano}
            & Cuauhtémoc   & 32.50 & 2.01 & 6.18 \\
            & Azcapotzalco & 33.50 & 2.01 & 5.99 \\
            & (resto)      &  ---  & 0.00 & 0.00 \\
        \hline
        \multirow{3}{*}{Mexicable}
            & Gustavo A. Madero &  87.84 & 2.80 & 3.19 \\
            & (resto)           &   ---  & 0.00 & 0.00 \\
            &                   &        &      &      \\
        \hline
        \multirow{3}{*}{Mexibús\textsuperscript{†}}
            & Iztacalco           & 23.08 & 1.52 & 6.58 \\
            & Venustiano Carranza & 33.84 & 1.90 & 5.62 \\
            & Gustavo A. Madero   & 87.84 & 3.23 & 3.67 \\
        \hline
    \end{tabularx}
    \vspace{0.4em}

    \noindent\textsuperscript{†}\,Datos de geometría y paradas integrados
    desde \termino{OpenStreetMap} \citep{osm2012} por ausencia de registro
    en el Portal de Datos Abiertos a la fecha de corrida.

    \fuentefigura{Elaboración propia con \texttt{VFTModel}
    \citep{cabrera2026VFTModel}; datos: \citep{adip_gtfs_2024, osm2012}.}
\end{table}

% ---------------------------------------------------------------
\subsubsection{Transporte de Superficie}
\label{subsubsec:cobertura-superficie}
% ---------------------------------------------------------------

Los sistemas que operan en vía compartida o mixta —\textsc{rtp}, Corredores
Concesionados, Trolebús, Pumabús y Tren Interurbano— presentan el patrón
territorial inverso: coberturas elevadas en las alcaldías centrales, pero con
una capilaridad que penetra en zonas periféricas donde la infraestructura
pesada no llega. La \textsc{rtp} y los Corredores Concesionados son los
sistemas con mayor cobertura absoluta de la red metropolitana, superando el
90\,\% en múltiples alcaldías.

\begin{table}[H]
    \centering
    \caption[Cobertura territorial — transporte de superficie, top 3 por alcaldía]{
    Tres alcaldías de la \cdmx{} con mayor cobertura territorial por sistema
    de transporte de superficie. Isócrono: 800~m. Corrida: 2026-04-18.}
    \label{tab:cobertura-superficie}
    \fontsize{8.5pt}{10.5pt}\selectfont
    \begin{tabularx}{\textwidth}{|l|X|r|r|r|}
        \hline
        \textbf{Sistema} & \textbf{Alcaldía} &
        \textbf{Área total (km²)} & \textbf{Área cub. (km²)} &
        \textbf{$C$ (\%)} \\
        \hline
        \multirow{3}{*}{\textsc{rtp}}
            & Miguel Hidalgo    & 46.39 & 41.90 & 90.32 \\
            & Cuauhtémoc        & 32.50 & 29.32 & 90.22 \\
            & Azcapotzalco      & 33.50 & 30.08 & 89.81 \\
        \hline
        \multirow{3}{*}{Corredores Conc.}
            & Cuauhtémoc          & 32.50 & 32.47 & 99.90 \\
            & Benito Juárez       & 26.68 & 26.32 & 98.65 \\
            & Iztacalco           & 23.08 & 22.25 & 96.42 \\
        \hline
        \multirow{3}{*}{Trolebús\textsuperscript{†}}
            & Cuauhtémoc    & 32.50 & 22.72 & 69.90 \\
            & Benito Juárez & 26.68 & 16.50 & 61.85 \\
            & Azcapotzalco  & 33.50 & 17.48 & 52.18 \\
        \hline
        \multirow{3}{*}{Pumabús}
            & Coyoacán      &  53.88 & 10.85 & 20.13 \\
            & Álvaro Obregón &  95.82 &  2.25 &  2.34 \\
            & Tlalpan        & 314.25 &  0.25 &  0.08 \\
        \hline
        \multirow{3}{*}{Tren Interurbano}
            & Cuajimalpa de Morelos &  71.10 & 1.59 & 2.23 \\
            & Álvaro Obregón        &  95.82 & 0.42 & 0.44 \\
            & (resto)               &   ---  & 0.00 & 0.00 \\
        \hline
    \end{tabularx}
    \vspace{0.4em}

    \noindent\textsuperscript{†}\,La Línea 13 del Trolebús no contaba con
    registro en el Portal de Datos Abiertos a la fecha de corrida; su
    geometría fue integrada desde \termino{OpenStreetMap}
    \citep{osm2012}. Los valores de cobertura excluyen esa línea.

    \fuentefigura{Elaboración propia con \texttt{VFTModel}
    \citep{cabrera2026VFTModel}; datos: \citep{adip_gtfs_2024, osm2012}.}
\end{table}

% ---------------------------------------------------------------
\subsubsection{Zonas de Mayor Déficit Territorial}
\label{subsubsec:cobertura-deficit}
% ---------------------------------------------------------------

El análisis de cobertura integrada —considerando la unión de los isócronos de
todos los sistemas— revela una inequidad territorial estructural: las
demarcaciones con menor acceso al transporte masivo son, de manera consistente,
las de mayor extensión territorial y menor densidad urbana consolidada, ubicadas
en las periferias sur y oriente de la cuenca.

\begin{table}[H]
    \centering
    \caption[Demarcaciones con menor cobertura de transporte masivo]{
    Cinco demarcaciones de la \cdmx{} y del Estado de México con menor
    cobertura territorial integrada del sistema multimodal. Isócrono: 800~m.
    Corrida: 2026-04-18.}
    \label{tab:cobertura-deficit}
    \begin{tabular}{|l|l|r|r|r|}
        \hline
        \textbf{Entidad} & \textbf{Demarcación} &
        \textbf{Área total (km²)} & \textbf{Área cub. (km²)} &
        \textbf{$C$ (\%)} \\
        \hline
        \multirow{5}{*}{\cdmx{}}
            & Milpa Alta              & 298.01 &  35.00 & 11.75 \\
            & La Magdalena Contreras  &  63.37 &  16.33 & 25.76 \\
            & Tlalpan                 & 314.25 &  94.31 & 30.01 \\
            & Tláhuac                 &  85.78 &  43.85 & 51.12 \\
            & Cuajimalpa de Morelos   &  71.10 &  37.12 & 52.20 \\
        \hline
        \multirow{5}{*}{Estado de México}
            & Tonanitla    &   8.52 & 0.00 & 0.00 \\
            & Capulhuac    &  21.50 & 0.00 & 0.00 \\
            & Calimaya     & 104.26 & 0.00 & 0.00 \\
            & Donato Guerra & 181.36 & 0.00 & 0.00 \\
            & Ayapango     &  50.83 & 0.00 & 0.00 \\
        \hline
    \end{tabular}
    \fuentefigura{Elaboración propia con \texttt{VFTModel}
    \citep{cabrera2026VFTModel}.}
\end{table}

La Figura~\ref{fig:grafica-cobertura-alcaldia} complementa esta lectura con
la distribución de cobertura integrada por alcaldía, permitiendo identificar
el gradiente centro-periferia que caracteriza la inequidad territorial de la
red metropolitana.

\begin{figure}[H]
    \centering
    \includegraphics[width=0.92\textwidth]{%
        Figures/Mapas/Tableau_Resources/RedActual/VFTModel/UniGrafica/%
        Grafica_VFTModel_2026_Cobertura_alcaldia}
    \caption[Cobertura territorial integrada por alcaldía de la \cdmx{}]{
    Cobertura territorial integrada de la red multimodal por alcaldía de la
    \cdmx{}, expresada como porcentaje del área cubierta por isócronos
    peatonales de 800~m. El gradiente centro-periferia es el patrón
    dominante.}
    \label{fig:grafica-cobertura-alcaldia}
    \fuentefigura{Elaboración propia con \textbf{Tableau};
    datos: \texttt{VFTModel} \citep{cabrera2026VFTModel}.}
\end{figure}

% ---------------------------------------------------------------
\subsubsection{Vista Cartográfica de la Cobertura Total}
\label{subsubsec:cobertura-cartografia}
% ---------------------------------------------------------------

Las Figuras~\ref{fig:mapa-cobertura-red-total} y~\ref{fig:mapa-cobertura-calor}
sintetizan la distribución espacial de la cobertura integrada mediante dos
representaciones cartográficas complementarias: el alcance territorial acumulado
de los isócronos de todos los sistemas y la densidad de esa cobertura como mapa
de calor. Ambas fueron elaboradas con \textbf{QGIS} a partir de los
\termino{GeoPackages} exportados por el \textbf{Modelo VFT}.

\begin{figure}[H]
    \centering
    \includegraphics[width=0.90\textwidth]{%
        Figures/Mapas/Qgis_Resources/RedActual/Mapa_Cobertura_Red_Total}
    \caption[Cobertura espacial total de la red multimodal — isócronas 800~m]{
    Cobertura espacial integrada de los once sistemas de transporte masivo de
    la \zmvm{}, calculada a partir de isócronos peatonales de 800~m. Las zonas
    sin cobertura se concentran en el sur y oriente de la cuenca metropolitana,
    evidenciando las demarcaciones con mayor exclusión territorial.}
    \label{fig:mapa-cobertura-red-total}
    \fuentefigura{Elaboración propia con \textbf{QGIS};
    datos: \texttt{VFTModel} \citep{cabrera2026VFTModel},
    GTFS \citep{adip_gtfs_2024} e INEGI 2023.}
\end{figure}

\begin{figure}[H]
    \centering
    \includegraphics[width=0.90\textwidth]{%
        Figures/Mapas/Qgis_Resources/RedActual/Mapa_Cobertura_Mapa_Calor}
    \caption[Densidad de cobertura de la red multimodal — mapa de calor]{
    Densidad de la cobertura espacial de la red multimodal expresada como mapa
    de calor. Los núcleos de máxima accesibilidad se ubican en las alcaldías
    centrales (Cuauhtémoc, Benito Juárez, Iztacalco), mientras que las
    periferias sur y oriente presentan densidades cercanas a cero, confirmando
    la inequidad territorial diagnosticada.}
    \label{fig:mapa-cobertura-calor}
    \fuentefigura{Elaboración propia con \textbf{QGIS};
    datos: \texttt{VFTModel} \citep{cabrera2026VFTModel},
    GTFS \citep{adip_gtfs_2024} e INEGI 2023.}
\end{figure}

Los datos de los once sistemas apuntan en la misma dirección: el acceso al
transporte masivo en la \zmvm{} no se distribuye de manera homogénea. Alcaldías
como Cuauhtémoc registran coberturas de hasta el 99.90\,\% bajo los Corredores
Concesionados; Milpa Alta, segunda demarcación en superficie de la \cdmx{}
solo superada por Tlalpan, no rebasa el 11.75\,\% en ningún sistema. Esa
distancia no expresa una diferencia operativa sino una acumulación histórica de
trazados e inversiones que reproduce, dentro del campo de la movilidad urbana,
la estratificación descrita por \citet{rendon2009ciudadanos}: ciudadanos con
acceso pleno a la red de oportunidades de la ciudad frente a quienes lo tienen
de forma parcial o ninguna. Sobre esta base territorial, la
Sección~\ref{subsec:fuerza-capilar} examina cómo esa escasez periférica se
concentra en un puñado de nodos que conectan la periferia con la red troncal.
